\chapter{Presupuesto y Planificación}
\label{chap:presupuesto}

Siguiendo la normativa de Trabajos Fin de Titulación de la ETSII -- UPM, este capítulo recoge el análisis económico estimativo del proyecto y el cronograma temporal de ejecución.

\section{Estimación presupuestaria}
\label{sec:presupuesto_costes}

El presupuesto global se desglosa en costes de hardware, herramientas de software y recursos humanos:

\begin{enumerate}
    \item \textbf{Costes de equipamiento hardware:}
    \begin{itemize}
        \item 3$\times$ Robots cuadrúpedos Luwu Dynamics XGO Mini 2S (incluyendo brazos robóticos articulados).
        \item 3$\times$ Sensores LiDAR 2D omnidireccionales y módulos Raspberry Pi Compute Module 5 Lite.
        \item Baterías Li-Ion auxiliares y fuentes de alimentación para banco de pruebas.
    \end{itemize}
    \item \textbf{Costes de software y licencias:}
    \begin{itemize}
        \item El 100\,\% de las herramientas software utilizadas (Ubuntu Linux, ROS 2 Jazzy, Docker Engine, BehaviorTree.CPP v4, Nav2 y compiladores GCC/Clang) se distribuyen bajo licencias libres y de código abierto (\textit{Open Source}), por lo que el coste directo en licencias es de 0,00\,\euro.
    \end{itemize}
    \item \textbf{Costes de personal (Recursos humanos):}
    \begin{itemize}
        \item Dedicación estimada de un Ingeniero de Desarrollo Robótico (horas dedicadas a arquitectura, programación en C++20, simulación, pruebas e integración).
    \end{itemize}
\end{enumerate}

\section{Planificación temporal y cronograma}
\label{sec:cronograma}

\textit{[Sección esquemática: Cronograma de fases y diagrama de Gantt que abarca la revisión bibliográfica, desarrollo del driver, diseño de Behavior Trees, simulación, pruebas en banco y redacción de la memoria].}
